\documentclass[12pt,a4paper]{article}

\usepackage[utf8]{inputenc}
\usepackage[T1]{fontenc}
\usepackage[margin=2.2cm]{geometry}
\usepackage{xcolor}
\usepackage{listings}
\usepackage{hyperref}
\usepackage{array}
\usepackage{longtable}
\usepackage{fancyhdr}
\usepackage{titlesec}

\hypersetup{
    colorlinks=true,
    linkcolor=blue,
    urlcolor=blue
}

\pagestyle{fancy}
\fancyhf{}
\lhead{SecureWatch}
\rhead{Source Code Explanation}
\cfoot{\thepage}

\titleformat{\section}{\Large\bfseries}{\thesection}{0.5em}{}
\titleformat{\subsection}{\large\bfseries}{\thesubsection}{0.5em}{}

\lstdefinestyle{bashstyle}{
    language=bash,
    basicstyle=\ttfamily\small,
    keywordstyle=\color{blue}\bfseries,
    commentstyle=\color{gray},
    stringstyle=\color{teal},
    breaklines=true,
    frame=single,
    columns=fullflexible,
    showstringspaces=false,
    tabsize=4
}

\title{
    \textbf{SecureWatch: Source Code Explanation}\\
    \large Bash Tool for SSH Brute Force Detection and Blocking
}
\author{Team Project}
\date{2026}

\begin{document}

\maketitle
\tableofcontents
\newpage

\section{Project Goal}

SecureWatch is a Bash security tool designed to detect repeated failed SSH login attempts.
The program reads an authentication log file, extracts IP addresses from lines containing
\texttt{Failed password}, counts the number of failed attempts per IP, compares the count
with a configurable threshold, and optionally blocks suspicious IP addresses using
\texttt{iptables}.

The tool was developed to satisfy the mini-project requirements: use Bash and standard
Linux commands, support several command-line options, implement process execution modes,
write logs, handle errors, and provide light, medium, and heavy tests.

\section{Global Architecture}

The project is organized into small Bash modules. Each file has a clear responsibility.

\begin{longtable}{|p{5cm}|p{9cm}|}
\hline
\textbf{File} & \textbf{Role} \\
\hline
\texttt{src/securewatch.sh} & Main controller. Reads options, loads configuration, chooses the action, and launches execution modes. \\
\hline
\texttt{src/config/config.conf} & Default configuration values such as threshold, log file, and output log paths. \\
\hline
\texttt{src/lib/parser.sh} & Reads the SSH log file and extracts IP addresses from failed login attempts. \\
\hline
\texttt{src/lib/detection.sh} & Counts attempts per IP and decides which IP addresses are suspicious. \\
\hline
\texttt{src/lib/blocking.sh} & Blocks suspicious IPs with \texttt{iptables} and restores blocked IPs. \\
\hline
\texttt{src/lib/logging.sh} & Writes events to \texttt{history.log} using the required format. \\
\hline
\texttt{src/lib/utils.sh} & Shared helper functions: colored output, file checks, root checks, and IP validation. \\
\hline
\texttt{scripts/install.sh} & Installs the command as \texttt{securewatch} so it can be executed from anywhere. \\
\hline
\texttt{tests/*.sh} & Light, medium, and heavy test scenarios. \\
\hline
\end{longtable}

\section{Main Program: \texttt{securewatch.sh}}

\subsection{Loading the Project Files}

The main script first detects its real location. This is important because the program can
be installed as a symbolic link in \texttt{/usr/local/bin/securewatch}. Even if the user runs
\texttt{securewatch} from another directory, the script can still find its configuration and
library files.

\begin{lstlisting}[style=bashstyle]
SCRIPT_PATH="${BASH_SOURCE[0]}"
while [[ -L "$SCRIPT_PATH" ]]; do
    SCRIPT_DIR="$(cd -P "$(dirname "$SCRIPT_PATH")" && pwd)"
    SCRIPT_PATH="$(readlink "$SCRIPT_PATH")"
    [[ "$SCRIPT_PATH" != /* ]] && SCRIPT_PATH="$SCRIPT_DIR/$SCRIPT_PATH"
done
SCRIPT_DIR="$(cd -P "$(dirname "$SCRIPT_PATH")" && pwd)"
\end{lstlisting}

After that, the script loads the configuration and all modules:

\begin{lstlisting}[style=bashstyle]
source "${SCRIPT_DIR}/config/config.conf"
source "${SCRIPT_DIR}/lib/utils.sh"
source "${SCRIPT_DIR}/lib/logging.sh"
source "${SCRIPT_DIR}/lib/parser.sh"
source "${SCRIPT_DIR}/lib/detection.sh"
source "${SCRIPT_DIR}/lib/blocking.sh"
\end{lstlisting}

\subsection{Command-Line Options}

The script uses \texttt{getopts} to parse Linux-style options.

\begin{longtable}{|p{2cm}|p{11cm}|}
\hline
\textbf{Option} & \textbf{Meaning} \\
\hline
\texttt{-h} & Display detailed help. \\
\hline
\texttt{-d} & Run detection only. \\
\hline
\texttt{-b} & Detect and block suspicious IPs. Requires administrator privileges. \\
\hline
\texttt{-r} & Restore/unblock previously blocked IPs. Requires administrator privileges. \\
\hline
\texttt{-l <dir>} & Store \texttt{history.log} and \texttt{blacklist.txt} in a custom directory. \\
\hline
\texttt{-s} & Run the action in a subshell. \\
\hline
\texttt{-f} & Run the action in the background without waiting. \\
\hline
\texttt{-t} & Run the action as a background job, then wait for completion. \\
\hline
\end{longtable}

\subsection{Action Selection}

The main script separates the action from the execution mode. For example, the action can
be detection, blocking, or restore, while the execution mode can be normal, subshell, fork,
or thread-style background job.

\begin{lstlisting}[style=bashstyle]
if [[ "$action" == "restore" ]]; then
    restore_blocked_ips
elif [[ "$action" == "block" ]]; then
    run_blocking
elif [[ "$action" == "detect" ]]; then
    run_detection
fi
\end{lstlisting}

\section{Execution Modes: \texttt{-s}, \texttt{-f}, and \texttt{-t}}

The mini-project asks for concepts related to subshell, fork, and thread execution. Bash does
not provide native threads like C or Java, so the project uses Bash background jobs to
simulate thread-style execution.

\begin{longtable}{|p{2cm}|p{4cm}|p{3cm}|p{5cm}|}
\hline
\textbf{Option} & \textbf{Code} & \textbf{Waits?} & \textbf{Explanation} \\
\hline
\texttt{-s} & \texttt{( execute\_action )} & Yes & Runs the action inside a child shell. \\
\hline
\texttt{-f} & \texttt{execute\_action \&} & No & Starts the action in the background and immediately returns. \\
\hline
\texttt{-t} & \texttt{execute\_action \& wait \$!} & Yes & Starts a background job and synchronizes with \texttt{wait}. \\
\hline
\end{longtable}

\begin{lstlisting}[style=bashstyle]
if [[ "$subshell" -eq 1 ]]; then
    log_event "INFOS" "Execution en mode subshell"
    ( execute_action )
elif [[ "$bg_fork" -eq 1 ]]; then
    execute_action &
    print_info "Execution en mode fork (PID: $!)"
    log_event "INFOS" "Execution en mode fork (PID: $!)"
elif [[ "$bg_thread" -eq 1 ]]; then
    execute_action &
    print_info "Execution en mode thread."
    log_event "INFOS" "Execution en mode thread"
    wait $!
else
    execute_action
fi
\end{lstlisting}

\section{Configuration: \texttt{config.conf}}

The configuration file defines default values. The user can override \texttt{LOG\_FILE} and
\texttt{THRESHOLD} from the command line.

\begin{lstlisting}[style=bashstyle]
THRESHOLD="${THRESHOLD:-5}"
LOG_FILE="${LOG_FILE:-/var/log/auth.log}"
HISTORY_LOG="${HISTORY_LOG:-logs/history.log}"
BLOCKED_IPS_LOG="${BLOCKED_IPS_LOG:-logs/blocked_ips.log}"
\end{lstlisting}

Example:

\begin{lstlisting}[style=bashstyle]
LOG_FILE=/tmp/auth-demo.log THRESHOLD=5 securewatch -d -l ~/securewatch-logs
\end{lstlisting}

\section{Parser Module: \texttt{parser.sh}}

The parser is responsible for reading the authentication log file and extracting IP addresses
from failed SSH login attempts.

\subsection{Input Validation}

Before reading the file, the module verifies that the log file exists and can be read.

\begin{lstlisting}[style=bashstyle]
if ! require_file "$log_file"; then
    log_event "ERROR" "Fichier de log introuvable : $log_file"
    return 102
fi

if [[ ! -r "$log_file" ]]; then
    print_error "Permission refusee pour lire : $log_file"
    log_event "ERROR" "Permission refusee pour lire : $log_file"
    return 103
fi
\end{lstlisting}

\subsection{IP Extraction}

The extraction uses standard Linux filters:
\texttt{grep} selects failed SSH attempts, and \texttt{grep -oP} extracts only the IP address.

\begin{lstlisting}[style=bashstyle]
grep "Failed password" "$log_file" 2>/dev/null | \
    grep -oP 'from \K[0-9]+\.[0-9]+\.[0-9]+\.[0-9]+' \
    > "$TEMP_IPS_FILE"
\end{lstlisting}

The extracted IP addresses are stored in:

\begin{lstlisting}[style=bashstyle]
/tmp/securewatch_ips.tmp
\end{lstlisting}

\section{Detection Module: \texttt{detection.sh}}

The detection module counts failed attempts per IP address and compares them with the
threshold.

\subsection{Counting Attempts}

The module uses a pipeline of Linux commands:

\begin{lstlisting}[style=bashstyle]
sort "$TEMP_IPS_FILE" | uniq -c | sort -rn
\end{lstlisting}

This pipeline sorts IP addresses, groups identical IPs, counts each group, and sorts the
result from highest count to lowest count.

\subsection{Suspicious IP Decision}

If an IP count is greater than or equal to the threshold, it is written to the suspects file.

\begin{lstlisting}[style=bashstyle]
if [[ "$count" -ge "$threshold" ]]; then
    echo "$ip:$count" >> "$TEMP_SUSPECTS_FILE"
fi
\end{lstlisting}

The suspicious IPs are stored in:

\begin{lstlisting}[style=bashstyle]
/tmp/securewatch_suspects.tmp
\end{lstlisting}

\section{Blocking Module: \texttt{blocking.sh}}

The blocking module uses \texttt{iptables}. Because firewall modification is a privileged
operation, the module checks that the script is running as root.

\subsection{Root Permission Check}

\begin{lstlisting}[style=bashstyle]
require_root || return 103
\end{lstlisting}

\subsection{Blocking an IP}

The script first checks whether the IP is already blocked. If not, it adds a DROP rule.

\begin{lstlisting}[style=bashstyle]
if iptables -C INPUT -s "$ip" -j DROP 2>/dev/null; then
    print_info "L'IP $ip est deja bloquee par iptables."
else
    iptables -A INPUT -s "$ip" -j DROP
fi
\end{lstlisting}

The IP is also saved in the blacklist file so it can be restored later.

\subsection{Restoring Blocked IPs}

The restore action reads the blacklist and removes each matching \texttt{iptables} rule.

\begin{lstlisting}[style=bashstyle]
iptables -D INPUT -s "$ip" -j DROP
\end{lstlisting}

After restore, the blacklist file is emptied.

\section{Logging Module: \texttt{logging.sh}}

The logging module writes messages to \texttt{history.log}. The required format is:

\begin{lstlisting}[style=bashstyle]
yyyy-mm-dd-hh-mm-ss : username : INFOS|ERROR : message
\end{lstlisting}

The timestamp and username are generated automatically:

\begin{lstlisting}[style=bashstyle]
timestamp="$(date '+%Y-%m-%d-%H-%M-%S')"
user_name="${USER:-$(whoami 2>/dev/null || printf 'unknown')}"
\end{lstlisting}

The function also normalizes old names such as \texttt{INFO} to \texttt{INFOS}, which helps keep
the log format compatible with the project requirement.

\begin{lstlisting}[style=bashstyle]
case "$type" in
    INFO)
        type="INFOS"
        ;;
    WARNING|WARN)
        type="INFOS"
        ;;
esac
\end{lstlisting}

\section{Utility Module: \texttt{utils.sh}}

The utility file contains common helper functions used by all modules.

\begin{itemize}
    \item \texttt{print\_info}, \texttt{print\_success}, \texttt{print\_warning}, and \texttt{print\_error} display colored terminal messages.
    \item \texttt{require\_file} checks whether a file exists.
    \item \texttt{require\_directory} checks whether a directory exists.
    \item \texttt{require\_root} verifies administrator privileges.
    \item \texttt{is\_valid\_ip} validates IPv4 addresses.
\end{itemize}

The IP validation first checks the general format, then verifies that every octet is between
0 and 255.

\section{Install Script: \texttt{scripts/install.sh}}

The install script prepares the project and creates a command named \texttt{securewatch}.
This allows the user to run the tool from any directory.

\begin{lstlisting}[style=bashstyle]
BIN_DIR="${INSTALL_BIN_DIR:-/usr/local/bin}"
COMMAND_NAME="${COMMAND_NAME:-securewatch}"
TARGET="$ROOT_DIR/src/securewatch.sh"
LINK_PATH="$BIN_DIR/$COMMAND_NAME"
ln -sf "$TARGET" "$LINK_PATH"
\end{lstlisting}

Typical installation on Linux:

\begin{lstlisting}[style=bashstyle]
cd ~/SSHGuard
bash scripts/install.sh
securewatch -h
\end{lstlisting}

\section{Testing Strategy}

The project includes three tests required by the assignment.

\begin{longtable}{|p{4cm}|p{4cm}|p{6cm}|}
\hline
\textbf{Test} & \textbf{Mode} & \textbf{Purpose} \\
\hline
\texttt{light\_test.sh} & Subshell \texttt{-s} & Small log. No IP should reach the threshold. \\
\hline
\texttt{medium\_test.sh} & Thread-style \texttt{-t} & One IP reaches the threshold and must be detected. \\
\hline
\texttt{heavy\_test.sh} & Fork \texttt{-f} & Large log with hundreds of failed attempts. Two IPs must be detected. \\
\hline
\end{longtable}

Run all tests:

\begin{lstlisting}[style=bashstyle]
bash tests/light_test.sh
bash tests/medium_test.sh
bash tests/heavy_test.sh
\end{lstlisting}

\section{Manual Demonstration Commands}

Create a fake authentication log:

\begin{lstlisting}[style=bashstyle]
cat > /tmp/auth-demo.log <<'EOF'
Jan  1 00:00:01 host sshd[1]: Failed password for invalid user admin from 192.168.1.10 port 111 ssh2
Jan  1 00:00:02 host sshd[2]: Failed password for invalid user admin from 192.168.1.10 port 112 ssh2
Jan  1 00:00:03 host sshd[3]: Failed password for invalid user admin from 192.168.1.10 port 113 ssh2
Jan  1 00:00:04 host sshd[4]: Failed password for invalid user admin from 192.168.1.10 port 114 ssh2
Jan  1 00:00:05 host sshd[5]: Failed password for invalid user admin from 192.168.1.10 port 115 ssh2
Jan  1 00:00:06 host sshd[6]: Failed password for invalid user test from 10.0.0.20 port 116 ssh2
EOF
\end{lstlisting}

Run detection:

\begin{lstlisting}[style=bashstyle]
LOG_FILE=/tmp/auth-demo.log THRESHOLD=5 securewatch -d -l ~/securewatch-demo-logs
\end{lstlisting}

Run the three execution modes:

\begin{lstlisting}[style=bashstyle]
LOG_FILE=/tmp/auth-demo.log THRESHOLD=5 securewatch -d -s -l ~/securewatch-demo-logs
LOG_FILE=/tmp/auth-demo.log THRESHOLD=5 securewatch -d -t -l ~/securewatch-demo-logs
LOG_FILE=/tmp/auth-demo.log THRESHOLD=5 securewatch -d -f -l ~/securewatch-demo-logs
\end{lstlisting}

Run blocking and restore:

\begin{lstlisting}[style=bashstyle]
sudo env LOG_FILE=/tmp/auth-demo.log THRESHOLD=5 securewatch -b -l /var/log/securewatch
sudo securewatch -r -l /var/log/securewatch
\end{lstlisting}

\section{Requirement Mapping}

\begin{longtable}{|p{6cm}|p{8cm}|}
\hline
\textbf{Requirement} & \textbf{Implementation} \\
\hline
Bash script for a real need & SSH brute force detection and blocking. \\
\hline
Use Linux commands and filters & Uses \texttt{grep}, \texttt{sort}, \texttt{uniq}, \texttt{awk}, \texttt{wc}, \texttt{iptables}, pipes, and redirection. \\
\hline
Mandatory options & Implements \texttt{-h}, \texttt{-f}, \texttt{-t}, \texttt{-s}, \texttt{-l}, \texttt{-r}; also adds \texttt{-d} and \texttt{-b}. \\
\hline
Administrator privileges & Blocking and restore require root privileges through \texttt{require\_root}. \\
\hline
Logging & Writes \texttt{history.log} entries with timestamp, username, type, and message. \\
\hline
Error handling & Uses error codes such as 100, 101, 102, and 103. \\
\hline
Light, medium, heavy tests & Implemented in \texttt{tests/light\_test.sh}, \texttt{tests/medium\_test.sh}, and \texttt{tests/heavy\_test.sh}. \\
\hline
Help documentation & Available with \texttt{securewatch -h}. \\
\hline
\end{longtable}

\section{Conclusion}

SecureWatch demonstrates a complete Bash-based security workflow. It combines log parsing,
regular expressions, file manipulation, pipes, process control, environment variables,
privilege checks, error management, and firewall commands. Its modular design makes the
source code easy to explain: the parser extracts IPs, the detector identifies suspicious
activity, the blocker handles firewall rules, the logger records events, and the main
controller connects everything through command-line options.

\end{document}
